% !TEX program = xelatex
\documentclass[11pt,a4paper]{ctexart}
\usepackage{amsmath,amssymb,amsthm}
\usepackage{geometry}
\geometry{margin=2.5cm}
\usepackage{booktabs}
\usepackage{hyperref}
\usepackage{xcolor}
\hypersetup{colorlinks=true,linkcolor=blue!60!black,urlcolor=blue!60!black}
\usepackage{enumitem}

\newtheorem{theorem}{定理}
\newtheorem{proposition}[theorem]{命题}
\newtheorem{lemma}[theorem]{引理}
\theoremstyle{remark}
\newtheorem*{remark}{注}

\newcommand{\diver}{\mathop{\mathrm{div}}}
\newcommand{\TT}{\mathbb{T}}
\newcommand{\QQ}{\mathbb{Q}}
\newcommand{\ZZ}{\mathbb{Z}}
\newcommand{\RR}{\mathbb{R}}
\newcommand{\CC}{\mathbb{C}}
\newcommand{\EE}{E_{11}}
\newcommand{\bN}{b_{11}}
\newcommand{\dd}{\,d}
\newcommand{\Isplit}{I_{\mathrm{split}}}
\newcommand{\FF}{\mathbb{F}}

\title{\textbf{Boyd 的 conductor 11 Mahler 测度猜想：研究报告}}
\author{科研文件夹 \texttt{boyd-conductor11/}}
\date{2026 年 8 月 4 日}

\begin{document}
\maketitle

\begin{abstract}
\noindent Boyd (1998) 猜想许多二元多项式的 Mahler 测度是椭圆曲线 $L$ 值的有理倍数。
conductor 11 的三条恒等式中两条已被 Brunault 证明，剩下一条（Boyd 1998 编号 (2-33)，
Samart 2023 eq.~(4.1)）仍开放。本报告：(i) 自包含地综述公开文献；
(ii) 将开放猜想数值确认到 \textbf{366 位}（$|\Isplit-\bN|=9.26\times10^{-367}$，PARI \texttt{lfun} 330 位交叉吻合；原公开记录 50 位）；
(iii) 发现并证明结构恒等式 $I_1+I_2=-m(S_0)$；
(iv) 用精确代数计算（非数值）证明 $x,y$ 是 $X_1(11)$ 上的 modular units（$5A=O$）；
(v) 第二波更正：第一波“边界非尖点阻断朴素 BMZ”的断言源于用错积分链——
正确的全圆带符号闭链 $\tilde\gamma$ 在 $H_1(E,\ZZ)^-$ 中拓扑闭合，
由此把猜想化为 $\int_{\tilde\gamma}\eta=2\pi b_{11}$，并给出
Beilinson--Brunault 路线的证明纲要；
(vi) 第三波：regulator 常数显式计算——除子、金刚石积、椭圆双对数全部精确/高
精度锁定，$D_E(P)=\frac{2\pi}{5}b_{11}$（40 位，精确吻合 Brunault (3.151)）；
早先笔记的因子 $2$（Touafek 转述）后经第八波文献核对证伪——Bloch 定理的
正确归一化在 conductor-17 曲线上是因子 $1$，k=0 的证明改锚 Bertin--Brunault 已证定理
（第十波：Brunault (3.210) 的 Weierstrass symbol 经 $\diamond$ 比率 $-1$ 桥接）；
(vii) 第四波：$S_k$ 族精确结构分析（环面交 $\Longleftrightarrow$ $k\in[-4,2]$、
模单位轨迹恰 $k=0,1$）+ conductor 53 机制失效判定（留作开放）；
(viii) 第五波：闭链引理严格化——$C'=\tilde\gamma+\beta_0=2\gamma^-$（比率先验
整数 + 15 位认证）；(ix) 第六波：整数识别的 Arb 球算术铁证化（比值球含 $-2$、
半径 $<1/2$），\eqref{C3} 由此\textbf{完全严格证毕}（\S\ref{sec:proofc3}）；
(x) 第十波：regulator 锚定的比率桥接重修（删除旧的循环论证）+ \eqref{C3} 推进到
\textbf{366 位}；
(xi) 第十一波（rev2）：比例引理锚定（引理 \ref{lem:ratio}）+ 反不变生成元本原性
（引理 \ref{lem:antiinv}）+ 符号的认证区间包围（\texttt{sign\_certify.py}，
$\Isplit\in[0.1489,0.1553]\subset(0,\infty)$）+ 认证计算定理（定理 \ref{thm:cert}）；
k=1 材料移入附录（\S\ref{sec:proofk1}），标注 conditional on 归一化引理。
\end{abstract}

\tableofcontents
\bigskip

\noindent\textbf{阅读指南}：第一部分（\S1--\S4）是公开文献的详细综述，自包含；
第二部分（\S5--\S8）是本次工作（数值攻击 + 代数分析）的结果，含第一波的更正、
证明纲要、regulator 常数计算（\S\ref{sec:regulator}）与 \eqref{C3} 的证明
（\S\ref{sec:proofc3}）。

\part*{第一部分：背景详解（公开部分）}
\addcontentsline{toc}{part}{第一部分：背景详解（公开部分）}

\section{Mahler 测度是什么}

\subsection{定义与 Jensen 公式}

对非零多项式 $P(x_1,\dots,x_n)\in\CC[x_1^{\pm1},\dots,x_n^{\pm1}]$，其\textbf{对数 Mahler 测度}
是 $\log|P|$ 在 $n$ 维环面 $\TT^n=\{|x_1|=\cdots=|x_n|=1\}$ 上的平均：
\[
m(P)=\int_0^1\!\!\cdots\!\int_0^1 \log\big|P\big(e^{2\pi i t_1},\dots,e^{2\pi i t_n}\big)\big|\dd t_1\cdots dt_n,
\qquad M(P)=e^{m(P)}.
\]

\textbf{单变量情形}完全可以算：若 $P(x)=a_0\prod_{j=1}^d(x-\alpha_j)$，Jensen 公式给出
\[
m(P)=\log|a_0|+\sum_{j=1}^d\log^+|\alpha_j|,\qquad \log^+v:=\max(\log v,0).
\]
即“首项系数 + 单位圆外根的贡献”。由此，整系数单变量多项式的 $m(P)$ 总是\textbf{代数数的对数}。
Kronecker 定理：$m(P)=0 \iff P$ 是分圆多项式乘以单项式。著名的 \textbf{Lehmer 问题}（1933）
问 $m(P)>0$ 能否任意小；目前最小纪录仍是 Lehmer 本人的
$m(x^{10}+x^9-x^7-x^6-x^5-x^4-x^3+x+1)=\log(1.17628\ldots)$。

\textbf{两变量情形}完全不同：$m(P)$ 一般是超越数，而且——这正是本课题的主题——
它竟然反复地等于\textbf{椭圆曲线 $L$ 函数特殊值的有理倍数}。

\subsection{降维：把二维积分化为一维}\label{sec:jensen}

计算 $m(P(x,y))$ 的实用方法是先对 $y$ 用 Jensen 公式。若把 $P$ 看成 $y$ 的多项式
$P=A(x)y^2+B(x)y+C(x)$，两根 $y_\pm(x)$，则
\[
m(P)=\frac1{2\pi}\int_0^{2\pi}\Big[\log|A(e^{i\theta})|
+\log^+|y_+(e^{i\theta})|+\log^+|y_-(e^{i\theta})|\Big]\dd\theta.
\]
本报告所有 Mahler 测度数值都是用这个一维积分算的（mpmath 任意精度）。

\subsection{第一个漂亮公式：Smyth (1981)}
\[
m(1+x+y)=\frac{3\sqrt3}{4\pi}L(\chi_{-3},2)=L'(\chi_{-3},-1)=0.3230659\ldots
\]
其中 $\chi_{-3}$ 是 conductor 3 的奇 Dirichlet 特征。证明路线：Jensen 降维后得到 log-sine 积分，
再认出它是 Clausen 函数（二重对数 $\mathrm{Cl}_2$）的特殊值，而
$\mathrm{Cl}_2(\pi/3)\propto L(\chi_{-3},2)$。
这个公式是后来一切“Mahler 测度 $=L$ 值”现象的鼻祖。注意曲线 $1+x+y=0$ 是\textbf{有理曲线}
（亏格 0），对应 Dirichlet $L$ 函数；亏格 1 时就轮到椭圆曲线的 $L$ 函数登场。

\section{椭圆曲线及其 $L$ 函数：conductor 11 为何特殊}

\subsection{椭圆曲线的 $L$ 函数与 conductor}

$\QQ$ 上的椭圆曲线 $E$ 有 $L$ 函数 $L(E,s)=\sum_{n\ge1}a_n n^{-s}$，其中好素数处
$a_p=p+1-\#E(\mathbb{F}_p)$。\textbf{conductor} $N$ 是衡量 $E$ 坏约化程度的正整数
（可看作“$E$ 的层级”）。\textbf{$\QQ$ 上椭圆曲线的最小可能 conductor 是 11}——
没有 conductor $1$ 到 $10$ 的椭圆曲线。所以“最简的”椭圆曲线 $L$ 值就是 conductor 11 的 $L$ 值，
这也是 Boyd 表格里 conductor 11 居首位的原因。

conductor 11 只有\textbf{一个同源类（isogeny class）}，含三条曲线（LMFDB \texttt{11.a1--a3}）。其中：
\begin{itemize}[nosep]
\item $X_1(11):\ y^2+y=x^3-x^2$（即 \texttt{11.a3}），判别式 $-11$，$j=-2^{12}/11$，
$E(\QQ)=\ZZ/5\ZZ$；
\item $X_0(11):\ y^2+y=x^3-x^2-10x-20$（即 \texttt{11.a1}）。
\end{itemize}

\subsection{模形式与函数方程}

模性定理（Wiles 等）：$L(E,s)$ 的系数 $a_n$ 是一个权 2、水平 $N$ 的尖形式 $f_E$ 的 Fourier 系数。
对 conductor 11，这个尖形式有极漂亮的 \textbf{eta 乘积}表达式：
\[
f_{11}(\tau)=\eta(\tau)^2\eta(11\tau)^2=q\prod_{n\ge1}(1-q^n)^2(1-q^{11n})^2=\sum_{n\ge1}a_nq^n,
\qquad q=e^{2\pi i\tau}.
\]

完备化 $L$ 函数 $\Lambda(E,s)=N^{s/2}(2\pi)^{-s}\Gamma(s)L(E,s)$ 满足函数方程
$\Lambda(E,s)=w\,\Lambda(E,2-s)$。$E_{11}$ 的根数 $w=+1$（秩 0）。由此推出本报告的核心常数恒等式：

\begin{proposition}
$b_N=L'(E,0)=\dfrac{N}{4\pi^2}L(E,2)$。
\end{proposition}
\begin{proof}
$s\to0$ 时 $\Gamma(s)=1/s+O(1)$；函数方程迫使 $L(E,0)=0$，故
$L(E,s)=L'(E,0)s+O(s^2)$，$\Gamma(s)L(E,s)\to L'(E,0)$，即 $\Lambda(E,0)=L'(E,0)$。
另一方面 $\Lambda(E,0)=\Lambda(E,2)=N(2\pi)^{-2}\Gamma(2)L(E,2)=\frac{N}{4\pi^2}L(E,2)$。
\end{proof}

所以 $L'(E,0)$（$s=0$ 处导数，Beilinson 猜想喜欢的点）与 $L(E,2)$（临界区间右端点，
模形式方法好算的点）只差一个有理因子，文献中两种写法混用。数值上
\[
\bN=L'(\EE,0)=\frac{11}{4\pi^2}L(\EE,2)=0.15214714172591804948622729747863\ldots
\]
（我们的 \texttt{code/b11.py} 用 $f_{11}$ 的系数和“近似函数方程”把它算到 800 位；
级数以 $e^{-2\pi n/\sqrt{11}}$ 速度收敛，尾项界驱动截断：442 项 @ 800 dps。）

\section{Boyd 猜想从哪来：Deninger、Beilinson 与 Boyd 的数值实验}

\subsection{Deninger 的观察（1995--1997）}

Deninger 研究 $K_2$ 与 Mahler 测度的联系时猜想
\[
m\Big(1+x+\frac1x+y+\frac1y\Big)\stackrel?=\frac{15}{4\pi^2}L(E_{15},2)=L'(E_{15},0),
\]
曲线 $1+x+1/x+y+1/y=0$ 恰为 conductor 15 的椭圆曲线。数值吻合到几十位。

\textbf{为什么应该有这种事？} 机制（Deninger--Rodr\'iguez Villegas 的解释）：
\begin{enumerate}[nosep]
\item 对“温顺”（tempered）多项式 $P$（Newton 多边形的每个面多项式都是分圆的，
等价于 $K$ 论中符号 $\{x,y\}$ 在曲线 $P=0$ 上 tame 平凡），Jensen 降维后的一维积分
可改写为\textbf{椭圆 regulator}
\[
\eta(x,y)=\log|x|\dd\arg y-\log|y|\dd\arg x
\]
沿环面与曲线相交路径的积分；
\item Beilinson 猜想（对模曲线已有定理级别的显式化）说这类 regulator 配对等于
$L(E,2)$ 的有理倍数。
\end{enumerate}
于是“$m(P)=r\,b_N$”是 Bloch--Beilinson 猜想的具体化身；Rodr\'iguez Villegas (1997)
进一步把 $m(P_k)$ 写成模形式，使得 CM（复乘）情形可以严格证明。

\subsection{Boyd 的系统实验（1998）}

Boyd（\emph{Experimental Mathematics} 7:1，引用 248 次）对形如
$P_k=A(x)y^2+(B(x)+kx)y+C(x)$ 的族做了大规模数值实验，按 conductor 分类列出几十条
$m(P)=r\,b_N$ 型猜想（$r$ 为小的有理数），每条验证到约 50 位小数。最小 conductor 就是 11。
此后 25 年，这些猜想被逐个击破：
\begin{itemize}[nosep]
\item CM 曲线（conductor 27, 32, 36）：Rodr\'iguez Villegas、Lal\'in--Rogers、Rogers--Zudilin 等；
\item conductor 14：Mellit (2012)、Mellit--Brunault（“平行线”椭圆双对数方法）、Touafek；
\item conductor 15：Rogers--Zudilin (2014) 证明 Deninger 原猜想；
\item conductor 20, 24：Rogers--Zudilin；conductor 21：Lal\'in--Samart--Zudilin (2016)；
\item \textbf{conductor 11：Brunault (2005/2006)}，见下节。
\end{itemize}

\subsection{Brunault 的突破与 BMZ 公式}\label{sec:bmz}

Beilinson 曾证明：模曲线上 modular units（除子只支撑在尖点上的有理函数）的 regulator
可以用 $L$ 值表示，但常数不显式。\textbf{Brunault 的博士论文（2005，ENS Lyon）把这个定理
完全显式化}，并应用于 $X_1(11)$——它恰好是可以用 modular units 参数化的椭圆曲线
（Brunault 后来证明这种曲线只有有限条）。由此他证明了 conductor 11 的两条 Boyd 猜想：
\begin{align}
m\big((1+x)(1+y)(1+x+y)+xy\big)&=\frac{77}{4\pi^2}L(\EE,2)=7\bN, \tag{C1}\label{C1}\\
m\big(y^2+(x^2+2x-1)y+x^3\big)&=5\bN. \tag{C2}\label{C2}
\end{align}

后来 Mellit--Brunault--Zudilin 把这类计算凝成一条可直接套用的公式
（\texttt{literature/zudilin-regulator.pdf}）：对 Siegel units
$g_a(\tau)=q^{NB_2(a/N)/2}\prod_{n\equiv a\,(N)}(1-q^n)\prod_{n\equiv -a\,(N)}(1-q^n)$，
\[
\int_{c/N}^{i\infty}\eta(g_a,g_b)=\frac1{4\pi}L(f_{a,b;c},2),
\]
其中 $f_{a,b;c}$ 是某个显式写出的权 2 模形式（Eisenstein 级数乘积组合）。
\textbf{直观含义}：只要多项式的 $x,y$ 是模曲线上的 modular units、且积分路径连接两个尖点，
regulator 积分就直接是一个 $L$ 值。这是目前证明此类恒等式的主力武器，
也是本报告 \S\ref{sec:proof} 分析证明路线时的标尺。

\section{仍然开放的 (C3)：torus 上有零点时会发生什么}\label{sec:c3}

\subsection{陈述}

Boyd 1998 编号 (2-33) 的族 $S_k=y^2+(x^2+kx+1)y+x^3$。取 $k=0$：
\[
S_0=y^2+(x^2+1)y+x^3=0\quad\Longleftrightarrow\quad \text{conductor 11 的椭圆曲线}.
\]

与 \eqref{C1}\eqref{C2} 的多项式不同，$S_0$ \textbf{在环面上有零点}：
$x=i$ 时 $y^2=i$，$x=-i$ 时 $y^2=-i$，即 $(x,y)=(i,\pm e^{i\pi/4})$ 与
$(-i,\pm e^{-i\pi/4})$ 满足 $|x|=|y|=1$。此时 Deninger 的 regulator 公式出现边界修正，
Boyd 数值发现
\[
m(S_0)=0.4056029559150104\ldots\ \ \text{“seemingly not } r\,\bN\text{”},
\]
$m(S_0)$ 与 $\bN$ 之间\textbf{未知任何有理关系}：我们用 PSLQ 在 $10^8$ 系数界内未发现（Boyd 亦报告了同样的阴性结果）。
但 Boyd 同时发现：\textbf{沿 branch cut 劈开的带符号积分}仍然等于 $\bN$。
写 $S_0=0$ 的两根
\[
y_\pm(x)=-\frac{x^2+1}{2}\pm\sqrt{\frac{(x^2+1)^2}{4}-x^3},
\]
$x=e^{i\theta}$ 沿上半环面走，$\theta=\pi/2$ 处正是 torus 交点 $x=i$（“分支切口”），定义
\[
\Isplit:=\underbrace{\frac1\pi\int_0^{\pi/2}\log|y_-(e^{i\theta})|\dd\theta}_{I_1}
\;-\;\underbrace{\frac1\pi\int_{\pi/2}^{\pi}\log|y_-(e^{i\theta})|\dd\theta}_{I_2}
\ \stackrel?=\ \pm \bN. \tag{C3}\label{C3}
\]
（符号取决于哪个根叫 $y_-$：$|y_+||y_-|=|x^3|=1$，换根整体变号。
Samart 的 eq.~(4.1) 把右端写成 $-L'(E,0)$；但按 Samart 自己的 $y_-$ 定义，左端积分是正数
$+0.1521471\ldots$，该负号疑似笔误——按我们的 $y_-$ 约定恒等式读作 $+\bN$，\eqref{C3} 因此记 $\pm$。）

\textbf{$y_-$ 的约定（澄清）}：$y_-$ 指 $|x|=1$ 上小模长根的\textbf{连续分支}：在 $\theta=0$ 处两根合并于 $y_-(1)=-1$；在 $(0,\pi)\setminus\{\pi/2\}$ 上由 $|y_-|<|y_+|$ 唯一确定；在第二个节点 $\theta=\pi/2$ 处按连续性延拓（两个节点处 $|y_\pm|=1$，故该处 $\log|y_-|=0$，被积函数不受取值选择的影响）。

Boyd 的原话（经 Samart 2023 \S4 引用）：
\begin{quote}
``This is in accord with our contention that in case $P$ vanishes on the torus, it is the integral
of $\omega$ around a branch cut rather than $m(P)$, which should be rationally related to $L'(E,0)$.''
\end{quote}

Samart 2023（arXiv:2301.05390）以 Boyd 的``$=$?''猜想记号把 \eqref{C3} 记录为\textbf{开放恒等式}
（其 eq.~(4.1)），并指出可尝试用他在
conductor 19 的 $Q_\alpha$ 族上成功的方法（超几何公式 + BMZ）来证，但“$S$ 族的情形 less apparent”。
这就是本次攻击的靶子。

\subsection{为什么这个情形难}\label{sec:whyhard}

\eqref{C1}\eqref{C2} 的证明链条是：modular units + 尖点间路径 + BMZ。
\eqref{C3} 的积分路径端点是 $\theta=0,\pi/2,\pi$ 对应的三个点，其中 $\theta=\pi/2$
（torus 交点）处曲线“穿过”环面；路径边界 $2[P_{\pi/2}]-[P_0]-[P_\pi]$ 是否由尖点组成、
能否套用 BMZ，文献中没有答案。第一波曾据此断言朴素 BMZ 被阻断；第二波发现该推理
用错了积分链——正确对象是全圆带符号闭链，它在 $H_1(E,\ZZ)^-$ 中自动闭合
（详见 \S\ref{sec:obstruction} 的更正与 \S\ref{sec:outline} 的证明纲要）。

\part*{第二部分：本次工作}
\addcontentsline{toc}{part}{第二部分：本次工作}

\section{方法}

\begin{itemize}
\item \textbf{$\bN$ 高精度}（\texttt{code/b11.py}）：由
$f_{11}=\eta(\tau)^2\eta(11\tau)^2=\sum a_nq^n$ 的整数系数
（Euler 函数平方的截断卷积，精确整数），用权 2、根数 $+1$ 的近似函数方程
\[
\bN=\Lambda(f,2)=\sum_{n\ge1}a_n\Big[e^{-t_n}\Big(\frac1{t_n}+\frac1{t_n^2}\Big)+E_1(t_n)\Big],
\quad t_n=\frac{2\pi n}{\sqrt{11}},
\]
$E_1$ 为指数积分（来自 $\int_1^\infty e^{-ty}/y\dd y$）。项衰减 $\sim e^{-1.894n}$；
尾项界驱动截断：442 项 @ 800 dps（旧稿``200 项够 300 位''有误，已更正）。
\item \textbf{Mahler 测度}：\S\ref{sec:jensen} 的一维 Jensen 积分，mpmath 80--300 dps。
\item \textbf{PSLQ}：\texttt{mpmath.pslq}，搜索小系数整数关系。
\item \textbf{椭圆曲线群律（精确）}：令 $u=2y+x^2+1$ 把 $S_0=0$ 化为四次曲线
$u^2=x^4-4x^3+2x^2+1$，在其上用首一抛物线 $u=x^2+bx+c$ 实现加法/取负
（\texttt{code/torsion.py}）；全部在 $\QQ$、$\QQ(\sqrt2)$、$\QQ(\zeta_8)$ 上用分数精确运算，
\textbf{不是数值近似}。
\end{itemize}

\section{数值结果}

\subsection{已证恒等式的独立复验（80 dps，\texttt{notes/attack1-results.txt}）}

\begin{center}
\begin{tabular}{@{}lll@{}}
\toprule
恒等式 & 计算值 & 与右端之差 \\
\midrule
\eqref{C1} & $1.06502999208142634\ldots$ & $|m-7\bN|\approx5.0\times10^{-53}$ \\
\eqref{C2} & $0.76073570862959024\ldots$ & $|m-5\bN|\approx3.6\times10^{-53}$ \\
\bottomrule
\end{tabular}
\end{center}

\subsection{开放猜想 \eqref{C3} 确认到 366 位 —— 主要数值结果（\texttt{code/attack13\_c3\_300.py}，存档 \texttt{notes/attack13-c3-300.txt}）}

\[
\Isplit=0.152147141725918049486227297478634495628143589164226122809889823882023289695302776676\ldots
\]
\[
|\Isplit-\bN|=9.26\times10^{-367}.
\]
（800 dps 工作精度，641 秒；$\bN$ 的 $\eta$-级数值与 PARI/GP \texttt{lfun(E,0,1)} \textbf{330 位逐位吻合}——三条独立计算共享 310 位公共前缀。原 \texttt{attack3.py} 的 149 位结果标记 superseded。）
此前公开记录是 Boyd 的 50 位验证（Boyd 2015 slides，p.~28，即 \texttt{literature/boyd-pnwnt2015.pdf}）；本次推进到 \textbf{366 位}。

\subsection{结构恒等式（先数值发现，后给出证明）}\label{sec:structure}

计算中注意到 $I_1+I_2=-m(S_0)$ 吻合到 152 位。事实上这是定理：

\begin{proposition}[区间算术证明]\label{prop:structure}
在 $[0,\pi]$ 上 $|y_-(e^{i\theta})|\le1\le|y_+(e^{i\theta})|$（球算术自适应二分认证：
$\log|y_-|$ 在 $(0,\pi)$ 上为负，仅在折点 $\theta=\pi/2$ 与端点 $\theta=0$ 处为零；
\texttt{code/branch\_certify.py}——本命题是 \S\ref{sec:proofc3} 意义上的计算机辅助结果）。
又 $|y_+y_-|=|x^3|=1$，故
\[
m(S_0)=\frac1\pi\int_0^\pi\log|y_+|\dd\theta=-\frac1\pi\int_0^\pi\log|y_-|\dd\theta=-(I_1+I_2).
\]
\end{proposition}

\begin{remark}
\eqref{C3} 因而等价于 $I_1=\dfrac{\bN-m(S_0)}{2}$、$I_2=-\dfrac{\bN+m(S_0)}{2}$。
也就是说，劈裂积分的猜想给出的是“大弧段积分”与“小弧段积分”各自的确切值。
\end{remark}

\subsection{$m(S_0)$ 的负结果}

$m(S_0)=0.40560295591501040\ldots$（Boyd 的 $0.4056029$ \checkmark）。
\begin{itemize}[nosep]
\item PSLQ$(m(S_0),\bN)$，系数界 $10^8$：\textbf{未发现关系}
（复核 Boyd 的 “seemingly not $r\bN$”）；
\item PSLQ 对 $\{m(S_0),\bN,\log2,\log3,\mathrm{Catalan},m(1+x+y)\}$，系数界 $10^{10}$：
\textbf{未发现关系}。支持“$m(S_0)$ 需用椭圆双对数表达而非初等常数”的预期。
\end{itemize}

\subsection{插曲：两个模型的 Mahler 测度不同}

Boyd slides 用模型 $P'=y^2+y+x^3+x^2$ 给出 $m=0.4056029$。我们算出
$m(P')=0.40560289185535\ldots$，而 $m(S_0)=0.40560295591501\ldots$——
\textbf{仅前 7 位相同}（差 $6.4\times10^{-8}$）。slides 的值是 $m(P')$；
这提醒“同一椭圆曲线的不同多项式模型，Mahler 测度不同”
（$m(P)$ 是多项式的不变量，不是曲线的不变量）。

\section{证明路线分析（上）：modular units 前提成立}\label{sec:proof}

劈裂积分可写成 regulator 积分：$|x|=1$ 上 $\log|y|\dd\arg x=-\eta(x,y)$，故
\[
\pi\,\Isplit=-\int_\gamma\eta(x,y_-),\qquad
\partial\gamma=2[P_{\pi/2}]-[P_0]-[P_\pi],
\]
\[
P_0=(1,-1),\quad P_\pi=(-1,-1+\sqrt2),\quad P_{\pi/2}=(i,e^{i\pi/4}).
\]

套 regulator 定理的前提：(i) $x,y$ 是 $X_1(11)$ 上的 modular units；
(ii) 积分链在 $H_1(E,\ZZ)^-$ 中闭合。

\textbf{(i) 成立（精确验证）}：四次模型的不变量给出 $j=-2^{12}/11=j(X_1(11))$ \checkmark。
群律精确计算（\texttt{code/torsion.py}，在四次模型上进行、以其无穷远点为群单位元）：对 $S_0$ 上的点 $A=(0,0)$（即四次坐标下的点 $(x,u)=(0,1)$），
\[
2A=(0,-1),\qquad 4A=(1,0)=-A\quad\Longrightarrow\quad \boxed{5A=O}.
\]
于是 $\diver(x)=[(0,0)]+[(0,-1)]-2[P_\infty]$ 与 $\diver(y)=3[(0,0)]-3[P_\infty]$
的支撑全是 5-扭点。这些恰为 $X_1(11)$ 的有理尖点：$S_0=0$ 与
$E:y^2+y=x^3-x^2=X_1(11)$ 的双有理识别是精确的（Riemann--Roch 计算复合 PARI 精确
极小化变换，往返验证，\texttt{code/kappa\_exact.py}），且在标准模同构（无穷远尖点
$\mapsto O$，$\omega_f=dx/(2y+1)$）下有理尖点 $P_v$（$v\in(\ZZ/11\ZZ)^\times/\pm1$）的像为
\[
P_1=\infty\mapsto O,\quad P_2\mapsto(1,0)=3A,\quad P_3\mapsto(0,-1)=4A,\quad
P_4\mapsto(0,0)=A,\quad P_5\mapsto(1,-1)=2A,
\]
即恰好映满 $E(\QQ)=\ZZ/5\ZZ$（Brunault Thm.~8 证明，(3.152)）。故除子支撑尖点化，
\textbf{$x,y$ 确为 modular units}（Manin--Drinfeld）。tempered 性同样精确化：Newton 面
多项式 $x^3+y$、$x^3+x^2y$、$x^2y+y^2$、$y^2+y$ 全分圆给出
$\{x,y\}\in K_2(E)\otimes\QQ$（Rodr\'iguez Villegas 判据）；更精确地，除子支撑处的
tame 符号由局部展开精确算出（\texttt{code/bertin\_diamond.py}）：
$T_v\{x,y\}=\pm1$（$v\in\{O,A,2A,3A\}$），故 $2\{x,y\}$ 的 tame 符号已平凡，
$\eta(x,y)$ 的实留数 $\pm\log|T_v\{x,y\}|$ 处处为零。

\section{闭性机制——第一波“障碍”分析的更正}\label{sec:obstruction}

\subsection{第一波的错误推理（如实记录）}

第一波曾在 $\QQ(\sqrt2)$、$\QQ(\zeta_8)$ 上算群律
（\texttt{code/endpoint\_torsion2.py}、\texttt{code/boundary\_torsion.py}），得到：
\begin{itemize}[nosep]
\item $P_0=(1,0)=-A$：5-扭点 \checkmark（尖点）；
\item $P_\pi=(-1,2\sqrt2)$、$P_{\pi/2}=(i,2\zeta_8)$：算到 $20P$ 均非扭点（坐标高度二次增长）；
\item 决定性组合 $T:=2P_{\pi/2}-P_\pi=(3,\,2i\sqrt2)$：算到 $30T$ 非扭点。
\end{itemize}
并据此断言：$\partial\gamma$ 不是尖点除子，朴素 BMZ（要求尖点间路径）被阻断。
\textbf{群律计算本身有效，但结论是错的}：它用错了积分链——把劈裂积分当成
“半圆上的路径积分”来取边界，而正确的对象是全圆上的带符号闭链。

\subsection{正确的积分链 $\tilde\gamma$ 是拓扑闭链}\label{sec:closed}

取全圆 $|x|=1$，$x=e^{i\theta}$，$\theta\in[-\pi,\pi]$，携带连续大模长分支
$y_+(\theta)$，权 $+1$（$\theta>0$）/ $-1$（$\theta<0$），在 $y_+$ 跨越单位圆的
折点 $\theta=\pm c$（$c=\pi/2$）处分段。这是 Samart（其 Lemma 9 的机制）意义下的
修正链 $\tilde\gamma$。

边界分析：链的边界只可能来自折点与端点。端点 $\theta=\pm\pi$ 给出
$[P_\pi]-[P_{-\pi}]$，而 $P_\pi=P_{-\pi}$（$x=-1$ 处两分支连续相接），相消；
折点贡献 $\partial\tilde\gamma=2\big([P_c]-[P_{-c}]\big)$。而
$P_{-c}=\overline{P_c}$（复共轭），故在 $H_1(E,\ZZ)^-$（复共轭的 $-1$ 特征空间）中
$[P_c]-[\overline{P_c}]$ 自动闭合——\textbf{闭性与 $P_c$ 是否为扭点无关}。
第一波找到的“障碍”（$T$ 非扭点）系假象：那是错误链的边界，正确链根本不经过那个组合。

\textbf{绕数计算}（60 位周期配对，\texttt{code/winding.py}、\texttt{code/winding.gp}）：
几何路径在折点 $\theta=\pm\pi/2$ 其实\textbf{不连续}（$y_{\mathrm{big}}$ 在两个交点
$(i,\pm e^{i\pi/4})$ 之间跳跃）——这正是必须取带符号链的原因；朴素环积分
$I_{\mathrm{loop}}=-0.47447\ldots i$ 甚至不是任何整闭链的周期（与 $w_{\mathrm{anti}}$
之比 $0.16262\ldots$ 非有理）。而带符号链在不变微分 $\omega=dx/u$
（四次模型 $u^2=x^4-4x^3+2x^2+1$）下的周期
\[
I_{\mathrm{signed}}=-2.917633233876990458\ldots i
\]
与 PARI 给出的 $H_1(E,\ZZ)^-$ 生成元周期 $w_{\mathrm{anti}}=2i\,\mathrm{Im}\,\omega_2$
（11.a3）\textbf{相等}（比值 $=1$ 到 13 位，受端点 $\sqrt\theta$ 奇性的积分误差所限）：
第二波据此断言``$\tilde\gamma$ 就是 $H_1(E,\ZZ)^-$ 的生成元（绕数 $n=1$）''。
\textbf{第五波更正与严格化}（\S\ref{sec:proofc3}）：$\tilde\gamma$ 是带边开链（端点
$P$ 非扭），其周期配对只是启发证据；正确的闭反不变闭链是
$C'=\tilde\gamma+\beta_0$（$\beta_0$ 为小分支补偿弧），严格认证
$\mathrm{class}(C')=2\gamma^-$——并由此把 $\int_{\tilde\gamma}\eta=2\pi\bN$
提升为\textbf{严格等式}。

\subsection{修正 Mahler 测度与 \eqref{C3} 的等价改写}

沿 $\tilde\gamma$ 的修正 Mahler 测度满足精确恒等式
\[
\tilde n=-\Isplit,\qquad \int_{\tilde\gamma}\eta(x,y)=2\pi\,\Isplit
\]
（\texttt{code/closedness\_check.py}，折点 $\theta=0,\pm\pi/2$ 分段，44 位）。于是
\[
\text{\eqref{C3}}\ \Longleftrightarrow\ \int_{\tilde\gamma}\eta(x,y)=2\pi\,\bN,
\]
数值上 $\tilde n=-\bN$ 到 44 位。

\subsection{$S_k$ 族：$\tilde n(k)$ 高精度数值表、PARI 独立鉴定与新恒等式}\label{sec:family}

把上述构造用于族 $S_k=y^2+(x^2+kx+1)y+x^3$（\texttt{code/ntilde\_family.py}，mpmath 50--80 位；
\texttt{code/verify\_family.gp}、\texttt{code/verify\_ratios.gp}，PARI/GP 2.15.5 独立复核）：
\begin{itemize}[nosep]
\item \textbf{折点存在 $\Longleftrightarrow$ $|k|<2$}，折角精确为 $c=\arccos(-k/2)$
（由 $2\cos\theta+k=0$，数值验证到 40 位）；$|k|\ge2$ 时 torus 无交点，闭链即全圆，$\tilde n(k)=m(k)$。
\item $E_k$ 的 $j$ 不变量（四次模型经典不变量 $I=k^4-8k^2+24k+16$、
$J=-2k^6+24k^4-72k^3-96k^2+288k-304$，$j=6912\,I^3/(4I^3-J^2)$）与 PARI \texttt{ellfromeqn}
全部一致；conductor 由 PARI 确认：
\[
k=-3,-2,-1,0,1,2,3\ \mapsto\ N=83,\;91=7\cdot13,\;53,\;11,\;17,\;37,\;79.
\]
\item $b$ 值双保险：我们的点计数管线（\texttt{code/b\_family.py}）与 PARI \texttt{lfun} 对到 50+ 位。
\end{itemize}

\begin{proposition}[环面交，精确刻画]\label{prop:torus}
对\textbf{实} $k$，记 $x=e^{i\theta}$、$y=e^{i\phi}$。环面上的点 $(x,y)$ 落在 $S_k=0$ 上当且仅当
\[
(k+2\cos\theta)\sin(\theta/2)=0\quad\text{且}\quad |k+2\cos\theta|\,|\cos(\theta/2)|\le2.
\]
由此：(i) 环面与曲线相交 $\Longleftrightarrow$ $k\in[-4,2]$；(ii) $-2\le k\le2$ 时出现\textbf{折点}
$\theta=\pm\arccos(-k/2)$（该处 $B=x^2+kx+1=0$，$y^2=-x^3$ 的两根模长均为 1）；
(iii) $-4\le k\le0$ 时另有 $\theta=0$ 处的\textbf{分支交换交点}（$S_k(1,y)=y^2+(k+2)y+1$
的两根都在单位圆上——$-4<k<0$ 时相异，$k=-4$（$y=1$）与 $k=0$（$y=-1$）处为重根即切触）；
(iv) $k=2$ 时折点与边界 $\theta=\pi$ 合并（$S_2(-1,y)=y^2-1$）。这与 Samart 观察到的
环面相交参数区间 $K\cap\RR=[-4,2]$ 精确吻合。
\end{proposition}

\begin{proof}[证明概要]
由 $x+x^{-1}=2\cos\theta$ 得 $B=x(k+2\cos\theta)$；方程
$e^{2i\phi}+Be^{i\phi}+e^{3i\theta}=0$ 除以 $e^{i\phi}$ 给出
$2\cos\psi\,e^{3i\theta/2}=-(k+2\cos\theta)\,e^{i\theta}$（$\psi:=\phi-3\theta/2\in\RR$）；
乘以 $e^{-i\theta}$ 后左端 $2\cos\psi\,e^{i\theta/2}$ 须等于实数 $-(k+2\cos\theta)$，
遂得两条件；反之满足两条件的 $\theta$ 给出 $\cos\psi\in[-1,1]$，从而给出 $\phi$。
在 $\theta\in[-\pi,\pi]$ 上 $\sin(\theta/2)=0$ 仅当 $\theta=0$（条件化为 $|k+2|\le2$，得
(iii)）；因子 $k+2\cos\theta=0$ 可解 $\Longleftrightarrow$ $|k|\le2$（得 (ii)）；两区间的并
为 $[-4,2]$，退化情形直接读出。
\end{proof}

\begin{center}\small
\begin{tabular}{cccccc}
\hline
$k$ & $N$ & $w$ & $c/\pi$ & $\tilde n(k)$ vs $b_N=|L'(E_k,0)|$ & 证据标签\\
\hline
$-3$ & $83$ & $-1$ & 无 & 比值 $0.8529175\ldots$ & 未发现 (PSLQ) \\
$-2$ & $91$ & $-1$ & 无 & 比值 $0.6339454\ldots$ & 未发现 (PSLQ) \\
$-1$ & $53$ & $-1$ & $1/3$ & 比值 $0.7392026\ldots$ & 机制失效（命题，见下）\\
$0$ & $11$ & $+1$ & $1/2$ & $\tilde n=-b_{11}$ & \textbf{已证}，\eqref{C3}，366 位 \\
$1$ & $17$ & $+1$ & $2/3$ & $\tilde n=+b_{17}$ & \textbf{已证}，定理（附录 A）\\
$2$ & $37$ & $-1$ & 无 & $m=\tilde n=2\,b_{37}$ & 数值，70 位 \\
$3$ & $79$ & $-1$ & 无 & $m=\tilde n=b_{79}$ & 数值，70 位 \\
$-4$ & $37$ & $-1$ & 无（相切） & $m=\frac72\,b_{37}$ & 数值，预测，25 位 \\
$-5$ & $359$ & $-1$ & 无 & $m=\frac14\,b_{359}$ & 数值，预测，25 位 \\
$-6$ & $997$ & $-1$ & 无 & $m=\frac18\,b_{997}$ & 数值，预测，25 位 \\
\hline
\end{tabular}
\end{center}

（``数值'' 行是猜想级恒等式；``未发现'' 指系数界 $10^8$ 的 PSLQ 搜索无有理关系。）

\begin{itemize}[nosep]
\item \textbf{$k=1$（conductor 17）}：$\tilde n(1)=b_{17}$——Samart 所述 conductor 17
“(4.1) 类似猜想”的独立高精度确认（$y_-$ 约定下积分 $=-L'(E_{17},0)$，有理因子 $r=1$）。
\item \textbf{$k=2,3$（conductor 37、79，根数 $-1$）}：无 torus 交点，Mahler 测度本身满足
Boyd 型恒等式 $m(S_2)=2|L'(E_{37},0)|$、$m(S_3)=|L'(E_{79},0)|$，70 位（PARI \texttt{lfun}）。
本次意外收获。
\item \textbf{$k=-1$（conductor 53）：机制失效判定（第四波，\texttt{code/k53\_attack.py} + \texttt{k53.gp}；第十波精确化 \texttt{code/k53\_smith.py}）}：
Samart 2023 \S4 关于 conductor 53 的``类似猜想''只有一句话（无公式、无精度、无引用）。
既然待检验的精确陈述不在案，我们无法直接裁决它；我们能够且确实检验的是：本报告的
\textbf{机制}——环面上闭反不变闭链配对 modular-unit symbol——能否延伸到 $k=-1$。
答案是否定的，有三个相互独立的原因：
\begin{enumerate}[nosep]
\item \textbf{拓扑障碍（精确）}：$k=-1$ 时 $|x|=1$ 上的 torus 交点为 $\theta=0$（两根 $e^{\pm2\pi i/3}$，
分支在此\textbf{交换}）与 $\theta=\pm\pi/3$（$y=\pm1$）。以交点为断点、big/small
分支组合的\textbf{闭反不变链}空间为 $\ZZ\cdot(1,1,1,1)$：反不变对称化的边界矩阵是
$\pm1/0$ 整数矩阵，其在 $\QQ$ 上的核\textbf{恰为一维}，由精确有理线性代数算出
（\texttt{code/k53\_smith.py}，非系数盒搜索）；生成元 $(1,1,1,1)$ 的周期为 $0$——而且是
\textbf{精确}为零，非数值：$(1,1,1,1)$ 是两叶在全圆上的和，两分支
$u=2y+B=\pm\sqrt{B^2-4x^3}$，故 $1/u_++1/u_-\equiv0$ 逐点成立，$\sum dx/u$ 恒等于零
（连续根沿圆走\textbf{两圈}才闭合，该两圈回路由同一抵消而为零）。因此
\textbf{环面上不存在非平凡的反不变闭链}，$H_1(E,\ZZ)^-$ 生成元无法在 torus 上实现，
Boyd 型环面积分对象根本不存在。
\item \textbf{代数障碍（精确除子论证）}：函数 $x,y$ 不是 modular units。事实上
\[
\diver(x)=[P]+[-P]-2[O],\qquad \diver(y)=3[P]-3[O],
\]
其中 $P=(0,0)$ 是 Mordell--Weil \textbf{生成元}：\texttt{53.a1} 有理扭子群平凡、秩 1
（PARI \texttt{ellorder}$(P)=0$、\texttt{ellidentify}）。该曲线是其 conductor 的
strong Weil curve，$X_0(53)$ 的两个尖点映到 $E$ 的扭点（Manin--Drinfeld），而扭子群
平凡故都映到 $O$——于是 $E$ 上的 modular unit 除子只能支撑于 $\{O\}$，而
$\diver(x)$、$\diver(y)$ 的支撑含非扭点 $P$。Beilinson--Brunault 机制因此没有可作用的
symbol，任何闭链的 regulator 配对都没有理由是有理数 $\times\,b_{53}$。
\item \textbf{数值观察}：自然候选对象（半环 $L_1$，连续根一圈，开链）周期与
$w_{\mathrm{anti}}$ 之比 $0.5492906\ldots$ 在已算精度内未见整数性，regulator 积分与
$2\pi b_{53}$ 之比 $0.7392026\ldots$；系数界 $10^8$ 的 PSLQ 搜索未发现有理关系。
\end{enumerate}
这些不排除经其他途径得到的恒等式；被排除的是证明 \eqref{C3} 的那个特定机制。
在缺乏精确候选公式的情况下，conductor 53 情形\textbf{留作开放}。
\item \textbf{扭点对照表}（\texttt{code/kfamily\_torsion.gp}）揭示真正的分野：
$k=0$ 扭子群 $\ZZ/5$、$\operatorname{ord}(0,0)=5$（Boyd 型恒等式\textbf{已证}，\eqref{C3}）；
$k=1$ 扭子群 $\ZZ/4$、$\operatorname{ord}(0,0)=4$（\textbf{已证}，\S\ref{sec:proofk1}）；
$k=-3,-2,-1,2,3$ 扭子群平凡、$(0,0)$ 无限阶（\textbf{非模单位}，命题，见下）。
\item \textbf{修正后的结构二分——精确结构分析 + 数值证据（第四波）}：分野不是 $k$ 的符号，而是
\textbf{torus 交点结构}：
\begin{enumerate}[nosep]
\item \textbf{族中的模单位（精确）}：(i) $k=0$（$E=X_1(11)$，$E(\QQ)_{\mathrm{tors}}=\ZZ/5$）与
$k=1$（$\ZZ/4$）时 $x,y$ 的除子支撑于有理扭点，而这些恰是模识别下的有理尖点（$k=0$ 见
\S\ref{sec:proof}，$k=1$ 见 \S\ref{sec:proofk1}），故 $x,y$ 是 modular units；这两种情形
Boyd 型恒等式\textbf{已证}（\eqref{C3} 与 \S\ref{sec:proofk1} 定理）。(ii)
$k\in\{-3,-2,-1,2,3\}$ 时 $E_k$ 有理扭平凡（PARI \texttt{elltors}，精确），$(0,0)\ne O$
非扭；每条 $E_k$ 是其 conductor 的 strong Weil curve，$X_0(N)$ 的尖点映到 $E_k$ 的扭点
（Manin--Drinfeld）即映到 $O$，modular unit 的除子只能支撑于 $\{O\}$，而 $x=0$ 与 $S_k$
交于 $(0,0)$ 与 $(0,-1)$。故这些 $k$ 的 $x,y$ \textbf{不是} modular units，模单位 regulator
机制不适用；$k=-1$ 时闭链成分也失效（上条命题）。
\item \textbf{观察到的二分（数值证据，与环面交命题联合）}：整数 $k\notin(-4,2)$（无真正
torus 交点）时标准 Deninger 机制预期直接适用，$m(S_k)=r_k\,|L'(E_k,0)|$（$r_k$ 为小有理数）
在所有已算情形获数值确认——$k=2,3$ 已确认，$k=-4,-5,-6$ 为\textbf{先预测后命中}：
$m(S_{-4})=\frac72|L'(E_{37},0)|$、$m(S_{-5})=\frac14|L'(E_{359},0)|$、
$m(S_{-6})=\frac18|L'(E_{997},0)|$（25 位精确）；注意 $S_{-4}$ 与 $S_2$ 同为 conductor 37，
给出同一曲线的 Rodriguez-Villegas 型关系。整数 $k$ 且 $-4<k<2$ 时，$x,y$ 为模单位的情形
恰是两个已证情形 $k=0,1$；$k=-1,-2,-3$（扭平凡、$\theta=0$ 处分支交换交点）在 PSLQ 界
$10^8$ 内未发现有理关系。我们把这些记录为\textbf{猜想的证据，而非定理}。
\end{enumerate}
\end{itemize}


\section{证明纲要：Beilinson--Brunault 路线}\label{sec:outline}

\textbf{主定理}：$\Isplit=\bN$（\eqref{C3} 中的带符号劈裂积分；$\bN=L'(E_{11},0)$）。

证明结合精确代数与计算机辅助步骤；除以下清单外的一切步骤都是精确的（有理或符号算术）：
\begin{itemize}[nosep]
\item \textbf{模序与分支指派}：结构命题（\S\ref{sec:structure}，命题 \ref{prop:structure}）的模长排序
与闭链的端点分支指派，球算术证明（\S\ref{sec:proofc3}）；
\item \textbf{同调类} $\mathrm{class}(C')=2\gamma^-$：由先验整数
$\mathrm{period}(C')/w_{\mathrm{anti}}$ 经球算术证明（\S\ref{sec:proofc3}）；
\item \textbf{regulator 积分的符号}：由严格含于半直线的认证区间包围证明（\S\ref{sec:proofc3}）；
\item \textbf{高精度浮点求值仅作一致性核查}（如与 $\bN$ 的 366 位吻合），从不作为证明步骤。
\end{itemize}

约定：$\gamma^-$ 是反不变生成元引理（\S\ref{sec:proofc3}，引理 \ref{lem:antiinv}）的\textbf{本原}
生成元；$\eta(x,y)=\log|x|\,d\arg y-\log|y|\,d\arg x$；$\ZZ[E]^-$ 取
Lal\'in--Ramamonjisoa 的商约定；$D_E$ 是背景部分的椭圆双对数，按 L--R Def.~5 eq.~(10) 逐字实现。

\begin{center}\footnotesize
\begin{tabular}{clll}
\hline
步骤 & 内容 & 状态 \\
\hline
S1 & 闭性引理：$C'=\tilde\gamma+\beta_0\in H_1(E,\ZZ)^-$，$\mathrm{class}(C')=2\gamma^-$（\S\ref{sec:proofc3}） & \textbf{已证（认证）} \\
S2 & tempered：$S_0$ 的 Newton 面多项式 & \\
   & $x^3+y$、$x^3+x^2y$、$x^2y+y^2$、$y^2+y$ 全分圆，故 $\{x,y\}\in K_2(E)\otimes\QQ$ & 已查 \\
S3 & modular units：$x,y$ 的除子支撑于尖点（$5A=O$ 精确验证，\S\ref{sec:proof}） & 已证 \\
S4 & Beilinson--Brunault 定理 + Brunault (3.151) & \\
   & $L(E,2)=\frac{10\pi}{11}D_E(P)$，系数 40 位复核 & 已核（文献+数值） \\
S5 & regulator 常数：Bertin--Brunault (3.210) 锚定 Weierstrass symbol $\{x_W,y_W\}$，经 $\diamond$ 比率 $-1$ 桥接（常数只依赖 $(E,\gamma^-)$）得 $\int_{\gamma^-}\eta=\pm2\pi\bN$（已证定理，\S\ref{sec:regulator} 第十波重修） & 闭合 \\
\hline
\end{tabular}
\end{center}

\subsection{regulator 常数的显式计算（第三波，40 位）}\label{sec:regulator}

第三波把 S5 的``余下代数计算''完成了（代码 \texttt{code/dilog.gp} + \texttt{code/dilog.py}，
原始输出 \texttt{notes/attack7-dilog.txt}）。全部成分如下。

\textbf{除子代数}（三次模型上的射影计算，Abel 主除子检验通过；\textbf{更正}：
此前误写 $\operatorname{div}(y)=3[A]-2[O]-[3A]$——$S_0$ 的射影闭包有\textbf{两个}
无穷远点 $O_1=[0:1:0]$、$Q_\infty=[1:-1:0]$（群单位元是 $Q_\infty$），分别映到
$3A$ 与 $O$，旧写法把二者混淆；$y$ 在 $A$ 处取值 $-1\ne0$，三重零点其实在 $2A$）：
\begin{equation}\label{eq:divs}
\operatorname{div}(x)=[A]+[2A]-[O]-[3A],\qquad \operatorname{div}(y)=3[2A]-2[3A]-[O],
\end{equation}
其中 $A=(0,0)$ 即 Brunault 记号中的 5-扭点 $P$，$O=[0:1:0]$ 为 11.a3 的群单位元。
金刚石积（修正后）
\begin{equation}\label{eq:diamond}
(x)\diamond(y)=6(O)-5(A)+5(2A),\qquad\text{故}\quad
D_E\bigl((x)\diamond(y)\bigr)=-5D_E(P)+5D_E(2P).
\end{equation}

\textbf{椭圆双对数}（mpmath 60 位；$\Delta<0$ 故取格基 $(w_1,\,w_1-w_2)$ 使
$\tau=0.5+0.2299i$ 落在上半平面，$q=e^{2\pi i\tau}\approx-0.2359$ 为负实数；
$z(P)=0.6w_1$、$z(2P)=0.2w_1$ 由 PARI \texttt{ellpointtoz} 给出）：
\[
D_E(P)=0.1911937370843316957549544343121738161012\ldots
\]
两条关键恒等式（均 40 位）：
\begin{itemize}[nosep]
\item \textbf{Brunault 系数}：$D_E(P)=\dfrac{2\pi}{5}\,\bN=\dfrac{11}{10\pi}L(E,2)$，
即 Brunault 论文 (3.151) 中系数的精确形式 $L(E,2)=\frac{10\pi}{11}D_E(P)$，数值钉死；
\item \textbf{exotic relation}：$D_E(2P)=\frac32\,D_E(P)$（Bertin 已证此类关系；
注意是 $3/2$ 而非朴素的 2 倍）。
\end{itemize}

\textbf{闭环}：代入 \eqref{eq:diamond} 与 exotic relation 得
\[
-5D_E(P)+5D_E(2P)=-5\Bigl(1-\frac32\Bigr)D_E(P)=\frac52\,D_E(P)=\pi\,\bN,
\]
而绕数一节已锁定 $\int_{\tilde\gamma}\eta(x,y)=2\pi\bN$（60 位，PARI 交叉验证）。

\textbf{regulator 的锚定（第十波重修；rev2 严格化）}：$\{x,y\}$ 沿 $\gamma^-$ 的 regulator
评估依赖两个输入：一个与归一化无关的\textbf{比例引理}，和一个把文献中被锚定的 symbol
识别出来的显式\textbf{字典}。

\begin{lemma}[比例引理]\label{lem:ratio}
固定 $E/\QQ$、$H_1(E,\ZZ)^-$ 的生成元 $\gamma$ 与 $\eta$、$\ZZ[E]^-$、$D_E$ 的约定。
设 Bloch 定理在这组数据上至多差一个约定常数：存在 $c\in\RR^\times$ 使对每个 tempered
symbol $\{f,g\}$，
\[
\int_\gamma\eta(f,g)=c\cdot D_E\big((f)\diamond(g)\big).
\]
则对任意两个 $\diamond$ 值非零的 tempered symbol，
\[
\frac{\displaystyle\int_\gamma\eta(f_1,g_1)}{\displaystyle\int_\gamma\eta(f_2,g_2)}
\;=\;\frac{D_E\big((f_1)\diamond(g_1)\big)}{D_E\big((f_2)\diamond(g_2)\big)}.
\]
\end{lemma}

\begin{proof}
假设的两端都是从 tempered symbol 群到 $(\RR,+)$ 的同态（regulator 侧由积分的线性，
$\diamond$ 侧由 $\diamond$ 的双线性与 $D_E$ 的可加性）；Bloch 定理即这两个同态的相等，
更换约定只把每端缩放一个固定非零标量（$\gamma$ 的符号、$D_E$ 的定向符号、$\ZZ[E]^-$
的商表示——后者下取值不变，因 $D_E$ 是奇函数），故比例假设即 Bloch 定理本身，$c$
吸收一切约定错配；特别地 $c$ 不依赖 symbol，两非零值相除即得结论。
\end{proof}

\begin{remark}
conductor-17 曲线（\S\ref{sec:proofk1}）上同一组约定以 $c=1$ 闭合，是有 L--R 佐证的
已证实例；本曲线上三个独立 tempered symbol 的认证积分恰给出 $c=2$（见下``因子 2''段）。
引理只用 $c$ 的存在性。
\end{remark}

\textbf{Brunault--Bertin symbol 字典}：文献中的锚定并不涉及我们的 symbol。Brunault
Thm.~8 固定 $E=X_1(11):y^2+y=x^3-x^2$ 及其 Weierstrass 坐标函数；(3.211) 定义
$r_\gamma\{f,g\}:=\frac1{2\pi}\int_\gamma\eta(f,g)$（$\gamma$ 生成
$H_1(E(\CC),\ZZ)^-$），(3.210)（引 Bertin 的 Crelle 版 Thm.~6 与 Cor.~6.1）评估该坐标对
的 $r_{\gamma^-}$。记该对为 $\{x_W,y_W\}$。Brunault 的 symbol 与 Weierstrass 坐标的识别
有三重独立保障：
\begin{enumerate}[nosep]
\item \emph{文本}：Brunault \S3.7--3.9 中名为 $x,y$ 的函数只有 Thm.~8 的坐标，且 (3.211)
是对 $E(\CC)$ 上的函数陈述的，而非对 Bertin 平面三次上的函数；
\item \emph{除子}：极小模型上
\[
\diver(x_W)=[A]+[4A]-2[O],\qquad \diver(y_W)=2[A]+[3A]-3[O],
\]
故在 $\ZZ[E]^-$ 中 $(x_W)\diamond(y_W)\equiv8(O)+5(A)-5(2A)$，由 Bertin exotic relation
得 $D_E(\diamond_W)=-\frac52D_E(P)$；
\item \emph{数值}：沿反不变生成元 $\gamma^-=s(2w_2-w_1)$（$0\le s\le1$）直接积分
$\eta(x_W,y_W)$，Richardson 外推 9 位得 $-2\pi\bN$（\texttt{code/bertin\_diamond.gp}），
与 (3.210) 吻合；Bertin 三次 $(X+1)(Y+1)(X+Y+1)+XY=0$ 的坐标 symbol 给出的却是
$+14\pi\bN$，由此排除。
\end{enumerate}

\textbf{应用到我们的 symbol}：比较 $\diamond$ 展开，
\[
D_E\big((x_W)\diamond(y_W)\big)=-D_E\big((x)\diamond(y)\big)\neq0
\]
（非零因 $D_E(P)=\frac{2\pi}{5}\bN\neq0$）。两个 symbol 都是 tempered 的：除子支撑处的
tame 符号皆为单位根（精确局部展开，\texttt{code/bertin\_diamond.py}：$\{x,y\}$ 在
$A,2A,3A,O$ 处 $T\in\{\pm1\}$；$\{x_W,y_W\}$ 有 $T_A=-1$、$T_O=+1$），故引理
\ref{lem:ratio} 适用并迫使
\[
\int_{\gamma^-}\eta(x,y)=-\int_{\gamma^-}\eta(x_W,y_W).
\]
Brunault (3.210) 给出
\[
|r_{\gamma^-}\{x_W,y_W\}|=\frac{5}{2\pi}D_E(P)=\bN,
\qquad\text{即}\qquad \int_{\gamma^-}\eta(x_W,y_W)=\pm2\pi\bN,
\]
符号取决于定向（Brunault Remarque~20：$D_E$ 依赖 $E(\RR)$ 的定向，只定义到符号）。故
\[
\int_{\gamma^-}\eta(x,y)=\pm2\pi\bN.
\]
\textbf{regulator 一侧由此成为已证定理}；旧稿``$K_2$ 秩 1 + 认证积分钉 $\lambda=1$''的
循环论证已删除。exotic relation $D_E(2P)=\frac32D_E(P)$
现同时引 Bertin 两篇：CRM Proc. Lecture Notes \textbf{36} (2004) 与
J.~Reine Angew. Math. \textbf{569} (2004) 175--188（后者即 Brunault 的参考文献 [10]）；
证明出处为 CRM 版（\textbf{更正}维持：Crelle 版中此关系仍是猜想；zbMATH 书评与
Touafek--Kerada、Mellit 三处佐证 CRM 版才是证明出处）。

\textbf{归一化问题（第十波定论）}：conductor-17 的证明中 Bloch 定理用
Lal\'in--Ramamonjisoa Thm.~6 的归一化：$r(\{x,y\})[\gamma]=D^E((x)\diamond(y))$、
$r[\gamma]=\int_\gamma\eta$。该因子-1 陈述\textbf{在该曲线上}的正确性由自洽性论证
保证（我们自己的重构；L--R 原文中该常数 $C$ 从未被确定）：L--R \S7 中实闭链的类是
$\gamma$ 的先验未知整数倍 $C$，将因子-1 定理用于其 $(X)\diamond(Y)$ 并与其已证的
Corollary~2（其 Thm.~1 eq.~(5)，结合 Zudilin 恒等式即其 eq.~(6)）比较，迫使
$C\cdot f=1$ 且 $C\in\ZZ\setminus\{0\}$，故 $f=1$；因子-2 陈述则迫使
$C=1/2\notin\ZZ$。其 \S5 的底层 $D_E$ 恒等式我们已 60 位复现。k=0 不用
$\diamond$-形式的绝对值：它只通过同一曲线上两个 tempered symbol 的\textbf{比率}
（上段，常数无论取何值都消去）进入，绝对值由 Bertin--Brunault 在 $\{x_W,y_W\}$ 上的
定理锚定。值得注意的是，因子-1 陈述本身在 conductor-11 曲线上\textbf{失效}：
那里检验的全部三个独立 tempered symbol 都满足
$\int_{\gamma^-}\eta=2\,D_E((\cdot)\diamond(\cdot))$（曲线级现象，见下段）；这不影响
任一证明——conductor-17 论证在自己的曲线上自洽，conductor-11 论证只用比率。
旧稿``$\pi r=D_E(\diamond)$（因子 2，Touafek 转述）''与两组数据都不符，已弃用。\textbf{因子 2 之谜的新数据（第十波，\texttt{code/bertin\_diamond.py} +
\texttt{bertin\_diamond.gp}，存档 \texttt{notes/attack13-bertin-diamond.txt}）}：
形式展开 $D_E((x)\diamond(y))=\frac52D_E(P)=\pi\bN$，而认证积分
$\int_{\gamma^-}\eta=2\pi\bN$——因子-1 归一化下两边恰差 2。我们用同一曲线上三个
\textbf{独立} tempered symbol 考察此事：Weierstrass symbol、我们的 $S_0$ symbol、以及
Bertin \eqref{C1} 三次 $(X+1)(Y+1)(X+Y+1)+XY=0$ 的坐标 symbol $\{X,Y\}$（经显式
Riemann--Roch 变换映到 $E$，PARI \texttt{ellidentify} 认证像为 \texttt{11.a3}，
精确变量代换 $[1,-1,-2,2]$）。$\diamond$ 值分别为 $-\frac52$、$+\frac52$、$+\frac{35}{2}$
倍 $D_E(P)$，而沿 $\gamma^-$ 直接积分 $\eta$ 给出 $-5$、$+5$、$+35$ 倍 $D_E(P)$
（前两个如上认证；第三个由外推 Riemann 和，与 $2\pi m(C_1)=14\pi\bN$ 一致）：
\textbf{因子在三个情形都恰为 2}。故它是 $(E,\gamma^-)$ 的曲线级性质，与 symbol、
$\diamond$ 约定（两曲线上相同）、tame symbol（皆单位根）、$D_E$ 级数无关。
conductor-17 的计算用同代码同约定闭合于因子 1（那里有 L--R 已证定理佐证）。
两条曲线恰在实拓扑上不同：$E$ 有 $\Delta=-11<0$、单实分支、$\Re\tau=1/2$
（故 $\gamma^-=a-2b$ 而非实轨迹），conductor-17 曲线有 $\Delta>0$、双实分支、
$\Re\tau=0$（故 $\gamma^-=b$）；\textbf{两种情形}下反不变生成元都本原
（引理 \ref{lem:antiinv} 与 17 情形的类似对合 $a\mapsto a$、$b\mapsto-b$），排除了
因子源于非本原闭链假象的可能。我们预期因子源于 Bloch 证明中单分支格的反不变闭链
对基本域边界的归一化，但未做机理级证明。由于 \eqref{C3} 的证明只用比率（常数消去），此问题不影响
主定理，如实记录。副产品：$D_E((X)\diamond(Y))=7\,D_E((x)\diamond(y))$ 精确成立，
在 $K_2(E)$ 层面解释了 Bertin $m(C_1)=7\bN$ 中的系数 7。

Samart 2023 仍将 \eqref{C3} 记录为开放恒等式，
缺口不在 regulator 计算，而在于 \textbf{Boyd 的劈裂积分链与 $H_1(E,\ZZ)^-$ 闭链
的等同}——这一环由第五波的 $C'=\tilde\gamma+\beta_0=2\gamma^-$ 认证闭合
（\S\ref{sec:proofc3}）。

\textbf{结论}：\eqref{C3} 的全部成分均已就位，且每一环要么是已证定理
（tempered、modular units、Bertin Thm 6、Brunault (3.151)/(3.210)、
Bloch Thm（k=1，Lal\'in 归一化）），
要么已被高精度数值 + 精确代数双重锁定（闭链 $C'=2\gamma^-$、除子 \eqref{eq:divs}、
$D_E$ 值）。\eqref{C3} 由此从``开放数值猜想''
降级为``已证定理的组装 + 书写级工作''。

\subsection{闭链引理的严格化与 \eqref{C3} 的证明（第五波）}\label{sec:proofc3}

完整细节见 \texttt{notes/proof-n1.md}；数值认证脚本与输出为
\texttt{code/n1\_certify.py}、\texttt{notes/attack9-n1.txt}。

\emph{约定（点值提升）}：每个环面分支 $\theta\mapsto y(e^{i\theta})$ 都默认视为点值提升
$\theta\mapsto(e^{i\theta},y(e^{i\theta}))\in E(\CC)$；因此下文
``$y_{\mathrm{big}}(c^-)=-P$'' 一类等式意指提升后的点在 $\theta=c^-$ 处等于 $E$ 上的点
$-P=(i,-e^{i\pi/4})$，而链的边界是 $E$ 上的除子。

\textbf{构造（精确代数）}。$\tilde\gamma$ 的边界为
$\partial\tilde\gamma=[P]-[\bar P]+[-P]-[-\bar P]=:D$，$P=(i,e^{i\pi/4})$
（精确：$x=i$ 时 $S_0=y^2-i$；分支跳跃值由数值读取并与精确值对照）。注意 $P$
\textbf{非扭}（第一波 \texttt{endpoint\_torsion2.py} 精确群律 20 倍无周期）——但闭性
不需要端点扭：取小分支补偿弧 $\beta_0=\alpha_1+\alpha_2$（$\alpha_2$：内弧
$P\to\bar P$；$\alpha_1$：外弧 $-P\to-\bar P$，内区间连接不了这两点是关键修正），则
$\partial\beta_0=-D$、$c(\beta_0)=-\beta_0$，于是
\[
C'=\tilde\gamma+\beta_0:\qquad \partial C'=0,\quad c(C')=-C',
\]
是闭的、反不变的整系数闭链。

\textbf{同调类（认证）}。属 1、$\Delta<0$ 时周期格有基 $(w_1,w_2)$：$w_1$ 实、
$\Re(w_2/w_1)=1/2$（单实分支）；对 \texttt{11.a3}，PARI \texttt{E.omega} 给出
$\tau=w_2/w_1=1/2+0.2298780212\ldots i$（上半平面约定）。反不变同调此时是显式的：

\begin{lemma}[反不变生成元]\label{lem:antiinv}
设 $E/\RR$ 有 $\Delta<0$、周期基 $(w_1,w_2)$ 如上，$(a,b)$ 为 $H_1(E,\ZZ)$ 的对偶基。
复共轭的作用为
\[
a\mapsto a,\qquad b\mapsto a-b.
\]
于是 $H_1(E,\ZZ)^-=\ZZ\,(a-2b)$，生成元 $a-2b$ \textbf{本原}，其周期
$\mathrm{period}(a-2b)=w_1-2w_2=-2i\,\mathrm{Im}\,w_2$ 是周期格唯一（差符号）的纯虚周期。
\end{lemma}

\begin{proof}
共轭固定实周期 $w_1$ 故固定类 $a$，把 $b$ 送到周期为 $\overline{w_2}=w_1-w_2$ 的闭链即类
$a-b$。类 $\alpha a+\beta b$ 反不变 $\Longleftrightarrow$
$(\alpha+\beta)a-\beta b=-\alpha a-\beta b$，即 $\beta=-2\alpha$，故
$H_1(E,\ZZ)^-=\ZZ(a-2b)$；$a-2b$ 本原因 $\gcd(1,2)=1$。其周期
$w_1-2w_2=-2i\,\mathrm{Im}\,w_2$ 纯虚；反之纯虚周期 $mw_1+nw_2$ 须满足 $2m+n=0$，
故为 $w_1-2w_2$ 的整数倍。
\end{proof}

周期配对 $\gamma\mapsto\mathrm{period}(\gamma)$ 把 $H_1(E,\ZZ)$ 等同于周期格，
故与 $\omega$ 的配对单射，$\mathrm{period}(\gamma^-)=\pm w_{\mathrm{anti}}$，于是
$\mathrm{period}(C')/w_{\mathrm{anti}}$ \textbf{先验为整数}。实测（60/80 位双精度对照）
\[
\mathrm{period}(C')=I_{\mathrm{signed}}+A_{s,\mathrm{outer}}-A_{s,\mathrm{inner}},\qquad
\frac{\mathrm{period}(C')}{w_{\mathrm{anti}}}=1.9999999999999999\ldots,
\]
$|\mathrm{period}(C')-2w_{\mathrm{anti}}|=2.5\times10^{-16}\ll 1$ ⟹
$\mathrm{class}(C')=2\gamma^-$（第六波 Arb 铁证：比值球含 $-2$、半径 $<1/2$，
符号=定向约定）。\textbf{更正}：开链 $\tilde\gamma$ 本身不闭，不能赋同调类；
此处所用的闭链 $C'$ 绕数为 \textbf{2}。

\textbf{同调不变性引理（与同调的配对）}：$\int_{C'}\eta(x,y)$ 只依赖于 $C'$ 在
$H_1(E,\ZZ)^-$ 中的同调类。事实上，三次模型 $y^2+y=x^3-x^2$ 的面多项式全分圆，由
Rodr\'iguez Villegas 的 tempered 判据，symbol $\{x,y\}$ 是 tempered 的：其全部 tame 符号
$T_v(\{x,y\})$ 皆为单位根（\texttt{code/bertin\_diamond.py} 精确验证，cf.\
\S\ref{sec:proof}）。于是 $\eta$ 在支撑每点 $v$ 处的\textbf{实}留数为
$\pm\log|T_v(\{x,y\})|=0$：$\eta$ 不带 delta 质量的留数，在支撑补集上为闭形式，在支撑点
处仅有（可积的）对数奇点。由此与 $\eta$ 的配对下降到同调，再经反不变性下降到
$H_1(E,\ZZ)^-$（引理 \ref{lem:antiinv}）。此处所用的链是 admissible 的：$C'$ 避开
$\operatorname{supp}\operatorname{div}(x)\cup\operatorname{supp}\operatorname{div}(y)$
（那些点 $x\in\{0,\infty\}$，而 $C'\subset\{|x|=1\}$）；折点处 $\log|y|=0$，故
$\eta=-\log|y|\dd\theta$ 沿 $C'$ 连续可积。

\textbf{regulator 合成（精确积分代数）}。$|x|=1$ 上
$\log|y_{\mathrm{big}}|=-\log|y_{\mathrm{small}}|$ 逐点成立（两根之积 $=x^3$），直接计算得
\[
\int_{\beta_0}\eta=2(J_1-J_2)=\int_{\tilde\gamma}\eta
\qquad\Longrightarrow\qquad \int_{C'}\eta=2\int_{\tilde\gamma}\eta,
\]
其中 $J_1=\int_0^{\pi/2}\log|y_s|d\theta$、$J_2=\int_{\pi/2}^{\pi}\log|y_s|d\theta$。
再由 \S\ref{sec:regulator} 第十波重修的锚定：Brunault (3.210)/(3.211) 给出
Weierstrass symbol $|r_{\gamma^-}\{x_W,y_W\}|=\frac{5}{2\pi}D_E(P)=\bN$
（$r=\frac1{2\pi}\int\eta$；并经 $\gamma^-=s(2w_2-w_1)$ 上的直接积分独立认证到
Richardson 9 位），两个 symbol 的 $\diamond$ 值比率恰为 $-1$、Bloch 常数只依赖
$(E,\gamma^-)$ 而消去，故 $\int_{\gamma^-}\eta(x,y)=\pm2\pi\bN$
（+ Brunault (3.151) $D_E(P)=\frac{2\pi}{5}\bN$
+ Bertin exotic $D_E(2P)=\frac32D_E(P)$，均为定理）。合成：
\[
\int_{\tilde\gamma}\eta=\frac12\int_{C'}\eta=\pm 2\pi\bN,
\]
故由结构恒等式 $\Isplit=\frac{1}{2\pi}\int_{\tilde\gamma}\eta$（\S\ref{sec:structure}，
已证）得 $|\Isplit|=\bN$。两个符号事实完成剩下的选择。\textbf{其一}，
$\bN=\frac{11}{4\pi^2}L(E,2)>0$：$L(E,s)$ 的 Euler 乘积在 $s=2$ 绝对收敛，且每个因子
$(1-a_p p^{-2}+p^{-1})^{-1}$ 为正，因由 Weil 界 $|a_p|\le2\sqrt p$ 有
$1-a_p p^{-2}+p^{-1}\ge1-2p^{-3/2}+p^{-1}>0$。\textbf{其二}，认证区间包围（球算术、
自适应二分加严格值域界；\texttt{code/sign\_certify.py}，证书
\texttt{notes/attack14-sign-k0.txt}）给出
\[
\Isplit\in[\,0.1489,\ 0.1553\,]\subset(0,\infty).
\]
故 $\Isplit=+\bN$：
\[
\boxed{\,\Isplit=\bN\,}\qquad\text{\eqref{C3} 证毕}.
\]

\textbf{区间算术铁证化（第六波完成）}：整数识别已由 Arb 球算术完全严格化
（\texttt{code/n1\_interval.py}，输出 \texttt{notes/attack10-interval.txt}）：
三段积分用 Arb 认证积分重算（$\theta=\pm t^2$ 换元消端点奇性 + Cauchy 尖端估计
——尖端常数 $a_0$ 由局部展开手算得到但不予轻信：脚本在 $t=\delta$ 用球算术重新评估 $f$，
四个尖端逐一认证 $|f(\delta)-a_0|\le H\delta^2$，$a_0$ 若有符号或分支错误会使值移动 $2$，
远超此界；
$D$ 穿负实轴处自适应细分 + 每段认证回避割线 + 认证符号传递；
$w_{\mathrm{anti}}$ 由 Newton+Rouch\'e 隔离的根 + Carlson RF 独立认证，与 PARI 45 位一致），
得比值球含 $-2$ 且半径 $<1/2$，先验整数性 $\Rightarrow$
$\mathrm{period}(C')/w_{\mathrm{anti}}=-2$ 为严格等式（符号=定向约定）。
模型常数 $\kappa=1$ 由判别式比 $\kappa^{12}=\Delta_{\mathrm{quartic}}/\Delta_{\min}=2^a11^b$
的离散候选 + 50 位一致锁定。

\textbf{两项实现注记}：(i) 存档输出（如 \texttt{notes/attack10-interval.txt}）中显示的
球半径遵循 Arb/python-flint 的 repr 约定——中点打印误差被折入显示半径；真实球半径更小
（显示区间包含真实球）。上文引用的所有认证界均用精确半径，而非显示半径。
(ii) 弧积分代码中，割线回避可能在某个子弧的首段选取 $i\sqrt{-D}$ 变体；现行脚本对
\textbf{每个}子弧（含首段）都认证叶选择。存档的 k=0 输出早于该逻辑的一次重构、已按其
复审：k=0 的割线穿越点位于子弧内部，其首段总用 $\sqrt D$ 变体，故认证值不受影响
（取错叶会使整段弧反号、比值以 $O(1)$ 幅度偏离 $-2$，存档输出排除了这一情形）。

\begin{theorem}[认证计算]\label{thm:cert}
以下各条由精确有理算术或认证区间（球）算术证明，全部证书存档于 \texttt{notes/}、
全部脚本在 \texttt{code/}（复现命令见 \S\ref{sec:reproduce}；软件：Python 3.12、mpmath、
sympy、python-flint 0.9.0/Arb、PARI/GP 2.15.5）：
\begin{enumerate}[nosep]
\item \textbf{精确代数输入}：从 $S_0=0$ 到极小三次 $E:y^2+y=x^3-x^2$ 的双有理映射与
$\omega_{\min}=dx/u$（$\kappa=1$）是精确多项式恒等式（\texttt{kappa\_exact.py}）；$x,y$
的除子、金刚石积、tame 符号、群律与扭子群 $E(\QQ)=\ZZ/5\ZZ$ 在
$\QQ,\QQ(\sqrt2),\QQ(\zeta_8)$ 上精确（\texttt{torsion.py}、\texttt{bertin\_diamond.py}）；
链端点 $P$ 的非扭性由好素数 $17$ 与 $89$ 处的既约证明（\texttt{endpoint\_torsion3.py}）。
\item \textbf{弧、定向、闭性}：\S\ref{sec:proofc3} 的链 $\tilde\gamma,\beta_0,C'$ 是
$|x|=1$ 上定向固定的显式参数化弧；端点分支指派（故 $\partial C'=0$、$c(C')=-C'$）在两种
尺度 $\varepsilon=10^{-6},10^{-9}$ 下经球算术认证（\texttt{branch\_certify.py}，证书
\texttt{notes/attack12-branch.txt}）。
\item \textbf{每段上的分支延拓}：周期积分的每个子弧上，所选平方根分支（$\sqrt D$ 或
$i\sqrt{-D}$）的解析性由整个子弧上的球求值认证——不限于端点附近——整体符号由认证的
匹配检验跨节点传递（\texttt{n1\_interval.py}）；结构命题的模序
$|y_{\mathrm{small}}|<1<|y_{\mathrm{big}}|$ 由自适应二分认证（\texttt{branch\_certify.py}）。
\item \textbf{区间包围}：每段弧积分与本原周期 $w_{\mathrm{anti}}$ 由 Arb 认证自适应积分与
Newton--Rouch\'e 根隔离的 Carlson $R_F$ 包围，积分半径 $\le10^{-3}$、周期半径
$\le10^{-48}$（\texttt{n1\_interval.py}，证书 \texttt{notes/attack10-interval.txt}）。
\item \textbf{整数识别}：类 $[C']\in H_1(E,\ZZ)^-=\ZZ\gamma^-$ 先验为整数（引理
\ref{lem:antiinv} 与同调类段）；认证比值球满足
$|\mathrm{ratio}-2|\le4.33\times10^{-3}<1/2$，最近整数唯一，故
$\mathrm{class}(C')=2\gamma^-$。
\item \textbf{符号}：最终符号由 $\Isplit$ 的认证包围固定：
$\Isplit\in[0.1489,0.1553]\subset(0,\infty)$（\texttt{sign\_certify.py}，证书
\texttt{notes/attack14-sign-k0.txt}），而 $\bN=\Lambda(f_{11},2)>0$ 由 $s=2$ 处绝对收敛的
Euler 乘积精确成立。
\end{enumerate}
主定理证明中唯一的非机器步骤是 \S\ref{sec:regulator} 的 regulator 评估，它使用 Brunault
的已发表定理（Thm.~8）与比例引理（引理 \ref{lem:ratio}）。
\end{theorem}

\section{总结}

\begin{enumerate}
\item \eqref{C1}\eqref{C2} 独立复现至 52 位；\eqref{C3} 确认至 \textbf{366 位}（$|\Isplit-\bN|=9.26\times10^{-367}$，PARI \texttt{lfun} 330 位交叉吻合；原公开记录 50 位，Boyd 2015 slides p.~28）。
\item 新结构定理 $|y_-|\le1\Rightarrow I_1+I_2=-m(S_0)$（命题 \ref{prop:structure}），
把 \eqref{C3} 化为 $I_1=(\bN-m(S_0))/2$。
\item $m(S_0)$ 对初等常数 PSLQ 阴性（界 $10^{10}$）。
\item modular units 前提\textbf{成立}（$5A=O$ 精确验证，\S\ref{sec:proof}）。
\item \textbf{更正}：第一波“朴素 BMZ 被非扭边界阻断”的断言不成立——正确积分链
$\tilde\gamma$ 在 $H_1(E,\ZZ)^-$ 中拓扑闭合，与扭点无关（\S\ref{sec:obstruction}）。
\item \textbf{\eqref{C3} 证毕（完全严格，\S\ref{sec:proofc3}）}：闭链引理严格化——
$\tilde\gamma$ 经小分支补偿弧闭化为 $C'=\tilde\gamma+\beta_0$（闭、反不变、整系数），
比率先验整数；15 位匹配认证 $\mathrm{class}(C')=2\gamma^-$（更正第二波``绕数 $n=1$''
的表述），并于第六波由 Arb 球算术铁证化（比值球含 $-2$、半径 $<1/2$）；配合
$\int_{\beta_0}\eta=\int_{\tilde\gamma}\eta$（精确积分代数）与
Bloch + Brunault (3.151)/(3.210) + Bertin exotic（均为定理；Brunault (3.210) 锚定的是 Weierstrass symbol $\{x_W,y_W\}$，经 $\diamond$ 比率 $-1$ 桥接到我们的 symbol，\S\ref{sec:regulator} 第十波重修），得
$\int_{\tilde\gamma}\eta=2\pi\bN$ 为\textbf{严格等式}，$I_{\mathrm{split}}=\bN$ 证毕。
\item 族结果（第四波完成 + rev2 精确化）：$\tilde n(k)$ 表 + \textbf{精确结构分析 + 数值证据}框架——
环面交精确刻画（命题 \ref{prop:torus}：交 $\Longleftrightarrow$ $k\in[-4,2]$）与模单位轨迹
精确刻画（恰 $k=0,1$）；$k\notin(-4,2)$ 时 Deninger 机制预期适用、$m(S_k)=r_k|L'(E_k,0)|$
获数值确认（$k=2,3$ 确认，$k=-4,-5,-6$ \textbf{先预测后命中}，$r=2,1,\frac72,\frac14,\frac18$）；
$-4<k<2$ 时恒等式证于 $k=0$（$\ZZ/5$）与 $k=1$（$\ZZ/4$），$k=-1,-2,-3$ PSLQ 阴性
（界 $10^8$ 内未发现）。\textbf{conductor 53 机制失效}：闭链空间精确一维、生成元周期由
$1/u_++1/u_-\equiv0$ 恒等于 0（环面上无反不变闭链可实现）+ $x,y$ 非 modular units
（精确除子论证：\texttt{53.a1} 为 strong Weil curve，$X_0(53)$ 尖点全映到 $O$，而
$\diver(x),\diver(y)$ 的支撑含非扭生成元）——机制不适用，在缺乏精确候选公式下
\textbf{留作开放}（\S\ref{sec:family}）。
\item \textbf{conductor 17（第七波，\S\ref{sec:proofk1}）}：同一方法再应用于
$k=1$——链结构逐字平行（$c=2\pi/3$、跳跃值 $y=\pm i$ 精确）、
$\mathrm{class}(C')=\pm2\gamma^-$ Arb 铁证、regulator 侧由 Lal\'in--Ramamonjisoa
已发表定理闭环——Samart 的 $\tilde n(1)=b_{17}$ 类比猜想成为定理，\textbf{conditional on
\S\ref{sec:proofk1} 讨论的归一化引理}。
\item \textbf{rev2（第十一波）严格化}：比例引理（引理 \ref{lem:ratio}，锚定常数与 symbol
无关）+ 反不变生成元本原性引理（引理 \ref{lem:antiinv}）+ 环面交/模单位的精确刻画
（命题 \ref{prop:torus}，\S\ref{sec:family}）+ 符号的认证区间包围
（\texttt{sign\_certify.py}、\texttt{k1\_sign\_certify.py}）+ 认证计算定理
（定理 \ref{thm:cert}）；k=1 材料移入附录并标注 conditional。
\end{enumerate}

\appendix

\section{k=1（conductor 17）：同一方法的再应用（第七波；conditional）}\label{sec:proofk1}

\textbf{本附录的地位}：正文在逻辑上独立于本附录。这里证明的定理（Samart 的 conductor-17
类比）使用 Lal\'in--Ramamonjisoa Thm.~6 的因子-1 归一化下的 Bloch 定理；``这是该曲线上的
正确归一化''来自下方归一化备注中的自洽性论证（我们自己的重构；L--R 原文并未确定该常数）。
因此本结果应读作 \textbf{conditional on 该归一化引理}；附录中其余一切与正文同一标准证明。

完整细节见 \texttt{notes/proof-k1.md}。Samart 的 conductor-17 类比猜想
$\tilde n(1)=b_{17}$ 沿 \S\ref{sec:proofc3} 的同一条路线\textbf{全程打通并证毕}，
认证级别与 k=0 相同：

\begin{itemize}[nosep]
\item \textbf{曲线}：$S_1=y^2+(x^2+x+1)y+x^3$，PARI \texttt{ellfromeqn} 给出
conductor \textbf{17}、$E(\QQ)_{\mathrm{tors}}=\ZZ/4\ZZ$、$(0,0)$ 为 4-扭
（平行于 k=0 的 5-扭）。关键差别：$\Delta=+17>0$，原初反不变周期直接是
$w_{\mathrm{anti}}=w_2$。
\item \textbf{闭链（精确代数）}：fold 角 $c=2\pi/3$；在 $\theta=c$，
$x=\omega=e^{2\pi i/3}$ 使 $x^2+x+1=0$，故 $S_1=y^2+1$，跳跃值\textbf{精确}为
$y=\pm i$；闭化构造逐字平行，$C'=\tilde\gamma+\beta_0$ 闭、反不变。
\item \textbf{同调类（Arb 铁证）}：\texttt{code/k1\_interval.py}
（真 300 位；$D(z)$ 在 $|z|=1$ 上无零点，$\min|D|=4$，无端点奇性；
$w_{\mathrm{anti}}$ 经 Newton+Rouch\'e + 两分支 Carlson RF 独立认证）：
比值球含 $-2$、半径 $3.4\times10^{-14}<1/2$ $\Rightarrow$
$\mathrm{class}(C')=\pm2\gamma^-$ 严格。输出 \texttt{notes/attack11-k1-interval.txt}。
\item \textbf{regulator（已发表定理闭环）}：除子形式与 k=0 全同
（局部展开 + 显式双有理映射 $X=-(x+y)$, $Y=x(x+y)$ + PARI + 代码展开四重验证）；
金刚石积 $(x)\diamond(y)=6(O)+4(A)-6(2A)$；$D_E(2A)=0$（2-扭，级数逐项为零，严格）；
$D_E(A)=\frac{17}{8\pi}L(E,2)$ 是 \textbf{Lal\'in--Ramamonjisoa 2017 已发表定理}
（\texttt{literature/lalin-ramamonjisoa-cond17.pdf}）；Bloch 用 Lal\'in Thm.~6
归一化（$\int_{\gamma^-}\eta=\pm D^E(\diamond)$，该曲线上有 L--R + Zudilin
已证恒等式背书）。
\item \textbf{合成}：$\int_{\tilde\gamma}\eta=\frac12\cdot2\cdot(\pm2\pi b_{17})$，
由结构恒等式得 $|\tilde n(1)|=b_{17}$；符号与主定理同样敲定：$b_{17}>0$ 由同一 Euler
乘积论证，认证区间包围（\texttt{code/k1\_sign\_certify.py}，证书
\texttt{notes/attack14-sign-k1.txt}）给出
$\frac1{2\pi}\int_{\tilde\gamma}\eta\in[-0.3026,-0.2961]\subset(-\infty,0)$，故
$\boxed{\tilde n(1)=+b_{17}}$（Samart conductor-17 类比猜想，证毕）。
\item \textbf{归一化备注（第八波定论）}：k=1 用的 Lal\'in 归一化
（$\int_{\gamma^-}\eta=\pm D^E(\diamond)$，因子 1）经 conductor-17 自洽性论证
确证为 Bloch 定理的正确形式（我们的重构：因子-1 定理与 L--R 已证 Corollary~2 比较
迫使 $C\cdot f=1$，$C\in\ZZ\setminus\{0\}$，详见 \S\ref{sec:regulator}；L--R 原文并未
确定 $C$）；k=0 的证明不用 $\diamond$-形式，直接锚定
Bertin--Brunault (3.210)--(3.211) 的已证定理（\S\ref{sec:regulator}）。
\end{itemize}

\section{复现方式}\label{sec:reproduce}

\begin{verbatim}
cd code && python b11.py && python attack1.py && python attack2.py \
  && python attack13_c3_300.py && python torsion.py && python endpoint_torsion2.py \
  && python boundary_torsion.py && python closedness_check.py \
  && python ntilde_family.py && python b_family.py && python winding.py \
  && python dilog.py && python k53_attack.py && python k53_smith.py && python kneg_m.py \
  && python n1_certify.py && python bertin_diamond.py \
  && python sign_certify.py && python k1_sign_certify.py
.venv/Scripts/python n1_interval.py    # Arb 区间算术铁证（需 python-flint）
.venv/Scripts/python branch_certify.py  # 分支指派 + 模序认证
.venv/Scripts/python k1_interval.py     # Arb 铁证，conductor 17
gp -q verify_family.gp && gp -q verify_ratios.gp
gp -q winding.gp && gp -q dilog.gp && gp -q bertin_diamond.gp
gp -q k53.gp && gp -q k53b.gp && gp -q kfamily_torsion.gp
gp -q k1_pari.gp && gp -q k1_points.gp && gp -q k1_zvals.gp
\end{verbatim}
依赖：Python 3.12 + mpmath + sympy；区间铁证另需 python-flint 0.9.0（项目内 \texttt{.venv}）。

\paragraph{Permanence（存档）}完整研究仓库——全部脚本、全部原始输出存档
（\texttt{notes/attack*.txt}）与本报告源文件——以 git 版本控制，并作为电子补充材料随
本报告发布；记录版本为标签 \texttt{rev2}（见仓库 log）。代码以 MIT license 发布。
上文 \texttt{attackN} 形式的文件名均指该存档。

\section{文献导读（\texttt{literature/}）}

\begin{itemize}[nosep]
\item \texttt{bertin-lalin-survey.pdf} — Bertin--Lal\'in 综述：全局图景与各 conductor 状态（先读这篇）
\item \texttt{boyd-pnwnt2015.pdf} — Boyd 2015 slides：猜想史 + $m(S_0)$ 原始数据；p.~28 载 \eqref{C3} 的 50 位验证（此前的公开纪录）
\item \texttt{brunault-these.pdf} — Brunault 博士论文：$X_1(11)$ 上 Beilinson 定理显式化，\eqref{C1} 的证明
\item \texttt{zudilin-regulator.pdf} — Zudilin：BMZ regulator 公式（证明武器）
\item \texttt{samart2023.pdf} — Samart：开放猜想 \eqref{C3} 的明确陈述（其 eq.~(4.1)）+ conductor 19 的成功范例
\item \texttt{lalin-samart-zudilin-cond21.pdf} — conductor 21：half-Mahler 方法范例
\item \texttt{lalin-ramamonjisoa-cond17.pdf} — Lal\'in--Ramamonjisoa 2017：conductor 17
已证公式 $L(E_{17},2)=\frac{8\pi}{17}D^E(P)$（k=1 证明的 regulator 闭环依据）与
Bloch Thm.~6 归一化出处
\item \texttt{boyd-slides.pdf} — Boyd 关于 $L(E,3)$ 的 slides
\end{itemize}
详细笔记：\texttt{notes/literature-notes.md}；原始运行输出：\texttt{notes/attack*-results.txt}。

\end{document}
